\scanpage{18}
% L1
expliciter les $H^{\unsure{q|n}}(C_{/\delta_i},F)$. Pour simplifier
% L2
la notation, on laisse tomber
% L3
l'indice $i$. \deleted{Notons que on} \inserted{Supposons}
% L4
\deleted{peut se \unsure{borner|besoins}}
% L5
que $\delta\in\mathcal N$ [i.e. $Z\in\operatorname{Ob}\,\underline{\unsure{S|\mathcal S|\mathfrak S}}$],
% L6
\deleted{donc} D'après un énoncé

% L3–6 left
\diagram{\begin{tikzcd}[ampersand replacement=\&,column sep=large,row sep=small]
X\arrow[r,"\delta"] \& Z.\\
T_0\arrow[u,"\varphi"]\arrow[r,"i"'] \& T\arrow[u,"u"']\\
T_0^{(i)}\arrow[u]\arrow[r] \& T^{(i)}\arrow[u,"u_i", "{\in\mathcal M}"'{draw,rounded corners=6pt,inner sep=1pt}]
\end{tikzcd}}
% L7
\unsure{standard|standart}, on peut remplacer $C_{/\delta}$ par
% L8
$(\mathcal M_{\unsure{C|0|{}}})_{/\delta}$, i.e. la catégorie des
% L9
diagrammes
\diagram{\begin{tikzcd}[ampersand replacement=\&,row sep=small,column sep=small]
X\arrow[r] \& Z\\
T_0\arrow[u,dash]\arrow[r,dash] \& T\arrow[u]
\end{tikzcd}}
qui sont des
% L10
\uline{morphismes \inserted{$\mathcal M$-}cartésiens} dans $C$, or cette
% L11
catégorie est \unsure{isomorphe|équivalente} à $\mathcal M_{/Z}$, avec
% L12
la topologie \unsure{héritée|héritie|induite} de $\unsure{\mathcal R|\mathcal T|\mathcal C}$. On trouve
% L13
donc
\[
H^q(C_{/\delta_i},F)\simeq H^q_{\mathcal M,\mathcal R}(\deleted{\unsure{X}}Z_n,F)
\]
% L14
et finalement
% L15
\[
\resizebox{.98\linewidth}{!}{$\displaystyle\boxed{0\longrightarrow\varprojlim\nolimits_n^{(1)}H^{\unsure{q|q-1}}_{(\mathcal M,\mathcal R)}(X_n^{(i)},F)
\longrightarrow H^q(X^{(i)},F)
\unsure{-|\to}\varprojlim\nolimits_n H^q_{(\mathcal M,\unsure{\mathcal T|\mathcal R|\mathcal C})}
\deleted{\unsure{(X_n}}(X_n^{(i)},F)\longrightarrow 0}$}
\]
% L16
[Sous l'hypothèse que les $X_n^{(i)}\in\operatorname{Ob}\,\underline S$]
% archive sheet number: 17

